% lemniro: {"kind":"note","description":"Open sets, neighborhoods, and the Hausdorff condition, with examples and a complete uniqueness proof.","topic":"Topology","series":"A first course in topology","order":1,"date":"2026-10-01","readingMinutes":12,"prerequisites":["Sets and functions","Unions and intersections"],"related":["compact-sets-are-closed","finite-intersections"]}
\documentclass[11pt]{article}
\input{tex/lemniro-preamble}
\title{Topological spaces and separation}
\begin{document}
\maketitle

\section{A language for nearness}
The usual distance on $\R$ tells us whether a point lies near another
point. Topology keeps the notion of an open neighborhood while allowing us
to forget the particular distance that produced it. The starting data are
a set and a distinguished collection of its subsets.

\begin{definition}[Topology]\label{def:topology}
A topology on a set $X$ is a collection $\T\subseteq\powerset{X}$ such that:
\begin{enumerate}
  \item $\varnothing\in\T$ and $X\in\T$;
  \item the union of any family of members of $\T$ belongs to $\T$;
  \item the intersection of any finite family of members of $\T$ belongs to $\T$.
\end{enumerate}
The pair $(X,\T)$ is a topological space, and the members of $\T$ are
called open sets. A subset $F\subseteq X$ is closed if $X\setminus F$ is open.
\end{definition}

The word ``finite'' in Definition~\ref{def:topology} is essential. For
example, each interval $(-1/n,1/n)$ is open in the usual topology on $\R$, but
\begin{equation}\label{eq:shrinking}
  \bigcap_{n\in\N}(-1/n,1/n)=\set{0},
\end{equation}
where $\N=\set{1,2,3,\ldots}$. The singleton on the right is not open: no
interval centered at zero fits inside it. This distinguishes a permitted
finite intersection from a countable intersection that need not be open.

\begin{example}[Three topologies]
The discrete topology on $X$ is $\powerset{X}$; every subset is open.
The indiscrete topology is $\set{\varnothing,X}$. On $\R$, the usual
topology consists of the sets $U$ for which every $x\in U$ has an
$\varepsilon>0$ satisfying $(x-\varepsilon,x+\varepsilon)\subseteq U$.
These examples share the axioms, but they contain very different
information about individual points.
\end{example}

\section{Neighborhoods and convergence}
An open neighborhood of $x\in X$ is an open set containing $x$.
In this note, whenever we say ``neighborhood'' without qualification we
mean an open neighborhood. A sequence $(x_n)_{n\in\N}$ converges to $x$ if
for every open neighborhood $U$ of $x$ there is an $N\in\N$ such that
$x_n\in U$ for all $n\geq N$.

This definition uses only open sets. In the usual topology on $\R$, it
agrees with the familiar epsilon definition: the intervals
$(x-\varepsilon,x+\varepsilon)$ suffice to test convergence because every
open neighborhood of $x$ contains one of them.

In an indiscrete space, every sequence converges to every point. There is
only one nonempty open set to test, namely the whole space. To exclude this
failure of uniqueness we impose a separation condition.

\section{The Hausdorff condition}
\begin{definition}[Hausdorff space]\label{def:hausdorff}
A space $X$ is Hausdorff if for every pair of distinct points $x,y\in X$
there are open sets $U,V\subseteq X$ such that
\[
  x\in U,\qquad y\in V,\qquad U\cap V=\varnothing.
\]
\end{definition}

\begin{example}
The usual real line is Hausdorff. If $x\ne y$, let
$r=|x-y|/3$. Then $(x-r,x+r)$ and $(y-r,y+r)$ are disjoint neighborhoods
of $x$ and $y$. Every discrete space is Hausdorff as well: use the
singletons $\set{x}$ and $\set{y}$. An indiscrete space with at least two
points is not Hausdorff, because its only nonempty open set is $X$.
\end{example}

\begin{proposition}\label{prop:unique-limit}
A sequence in a Hausdorff space has at most one limit.
\end{proposition}
\begin{proof}
Suppose that a sequence $(x_n)$ converges both to $x$ and to $y$, and
assume for a contradiction that $x\ne y$. By
Definition~\ref{def:hausdorff}, choose disjoint open neighborhoods $U$ of
$x$ and $V$ of $y$. Convergence to $x$ provides an index $N_U$ such that
$x_n\in U$ whenever $n\geq N_U$. Convergence to $y$ similarly provides
$N_V$ such that $x_n\in V$ whenever $n\geq N_V$. For
$n\geq\max\set{N_U,N_V}$, both conclusions hold, so
$x_n\in U\cap V=\varnothing$, which is impossible. Hence $x=y$.
\end{proof}

\begin{remark}
Proposition~\ref{prop:unique-limit} proves one consequence of the
Hausdorff condition. It does not say that uniqueness of sequence limits
characterizes Hausdorff spaces. Sequences alone need not detect all
topological behavior in an arbitrary space.
\end{remark}

\section{What to keep in view}
The axioms allow arbitrary unions but only finite intersections.
The Hausdorff condition lets us separate two points by disjoint open
neighborhoods. In the next note, compactness will turn infinitely many
pointwise separations into finitely many, and the finite-intersection
axiom will then produce a single open neighborhood.

\begin{exercise}
Show that every singleton is closed in a Hausdorff space. Fix $x\in X$,
and for each $y\ne x$ find an open neighborhood of $y$ that does not
contain $x$. Explain why the union of those neighborhoods is exactly
$X\setminus\set{x}$.
\end{exercise}
\end{document}
